\documentclass[10pt]{article}
\usepackage[a4paper,margin=0.85in]{geometry}
\usepackage{amsmath,amssymb,amsthm,booktabs,listings}
\usepackage[T1]{fontenc}
\usepackage[colorlinks=true,allcolors=blue]{hyperref}
\newcommand{\Inst}{\mathsf{Inst}}
\newcommand{\Norm}{\mathsf{N}}
\newcommand{\Capture}{\mathsf{Capture}}
\newcommand{\Sync}{\mathsf{Sync}}
\newcommand{\Eval}{\mathsf{Eval}}
\newcommand{\Req}{\mathsf{Require}}
\newcommand{\Ens}{\mathsf{Ensure}}
\newcommand{\FV}{\mathsf{FV}}
\newcommand{\sem}[1]{[\![#1]\!]}
\newtheorem{proposition}{Proposition}
\newtheorem{definition}{Definition}
\title{Instance States and Continuous Normalization in TeSpec\\
\large Operational semantics and implementation correspondence}
\author{Implementation note}
\date{October 10, 2026}
\begin{document}
\maketitle

\section{Problem and semantic scope}
TeSpec tests a partially instantiated separation-logic contract by symbolic
execution. It supports two control-flow policies: execution of loop iterations
and available function bodies (\texttt{expand}), and QCP's invariant and
contract abstraction (\texttt{no-expand}). Both policies maintain current
instance information at program points and evaluate known logical expressions.
The distinction concerns control-flow abstraction, not whether executable
logical definitions may be unfolded.

Let $\sigma=(\rho,h,\eta)$ consist of a C variable environment, a heap, and a
logical valuation. We use QCP's interpretation of typed storage, structure
layout, arithmetic, and separating conjunction. In $A*B$, the two assertions
hold on disjoint component heaps. In $A\land B$, they hold on the same heap.

A contract is $\{R(\bar a_0,\bar w)\}\ f(\bar a)\ \{E(\bar a_0,\bar w,r)\}$.
The tuple $\bar a_0$ denotes immutable entry snapshots and $\bar w$ denotes
contract logical parameters. A typed partial input substitution $\theta_0$
fixes selected snapshots and logical parameters. Pointer values may remain
symbolic while logical heap contents are fixed.

When a substitution occurs as a proposition, it denotes its graph:
$\bigwedge_{x\mapsto v\in\theta_0}x=v$.
Two queries must be distinguished:
\begin{align}
 &\exists\eta,h.\ (\eta,h)\models R[\theta_0], \label{eq:sat}\\
 &\{R[\theta_0]\}\ \mathit{body}\ \{E[\theta_0]\}. \label{eq:verification}
\end{align}
An existential witness for~\eqref{eq:sat} is not installed as an execution input
for~\eqref{eq:verification}. Unbound inputs in the latter range over all
interpretations satisfying the instantiated precondition.

\section{Product states}
\begin{definition}[Test state]
A branch at program point $\ell$ carries
\[
 T=\langle\ell,P,I,\theta_0,\Omega,m\rangle,
 \qquad I=\langle\nu,L,K,M\rangle.
\]
$P$ is the current QCP assertion, $I$ is the current instance view, $\Omega$ is
the set of outstanding obligations, and $m$ is the execution mode. $L$ maps
program-variable identities to symbolic storage locations, $K$ contains valid
logical equalities, and $M$ records typed storage and spatial predicates.
$\nu$ identifies an observation version.
\end{definition}
The instance view is branch-local and owns its data. Its heap component uses
QCP's existing spatial representation; it does not allocate numeric addresses
or introduce a second concrete heap semantics. Opaque predicates may remain
symbolic entries in $M$.

Write $\gamma(I)$ for the consistency relation represented by an instance
view. The conceptual product denotes
\[
 \sem{T}=\{\sigma\mid\sigma\models P\land\gamma(I)\land\theta_0\}.
\]
The implementation uses a materialized representation: valid instance
refinements are included in the assertion before a cut, and synchronized
views satisfy $P\models\gamma(I)$. In particular, the product is not encoded
as $P*\gamma(I)$, which would claim a second copy of the same storage.

Reconstructing $I$ from a weakened assertion alone loses information.
Consequently, an annotation cut first transports available instance bindings
to its target, generates the corresponding obligation, and only then captures
the successor view. Ordinary execution rebuilds the view from QCP's updated
formula; it never replays entry heap contents over a modified heap.

\section{Instantiation and evaluation}
\begin{definition}[Logical instantiation]
$F[\theta]$ is typed, capture-avoiding substitution by logical identity.
It removes instantiated existential binders and freshens independent binders
when needed. A symbolic pointer substituted for an existential is a witness
term, not a request to rename the existential binder to that pointer.
\end{definition}
Entry snapshots, current C values, and callee post-state variables have distinct
identities even when their source names coincide. Define point instantiation by
\[
 \Inst_T(F)=\Norm\bigl(\mathsf{Resolve}(F[\theta_0];L,K,M)\bigr).
\]
Resolution uses the current storage and valuation, with QCP's obligations for
memory reads and C operations. Normalization evaluates known arithmetic,
Boolean, list, and constructor operations, and unfolds accessible definitions
under a budget. Unknown subexpressions remain symbolic.

For a valid branch context $\Gamma$, the required normalization property is
\[
 \theta_0\land\Gamma\models F\leftrightarrow\Norm_\Gamma(F).
\]
Fresh values introduced by a definition expansion are existentially closed on
the right. C safety conditions remain obligations; mathematical simplification
must not erase overflow, division-by-zero, invalid-shift, or storage conditions.

For example, $n_0=2$ and $xs_0=[10,20]$ describe the entry instance. After
\texttt{n++; p[0]=n;}, the current value is $n=3$ and current contents are
$[3,20]$, while $xs_0$ remains $[10,20]$.

\section{Common operational rules}
For an ordinary statement $c$, the common transition is
\[
\frac{
 \widehat P=\mathsf{Materialize}(P,I,\theta_0)\quad
 (P_r,V)=\mathsf{QCPstep}(c,\widehat P)\quad
 P'=\Norm(P_r)\quad I'=\Capture(P')
}{
 \langle\ell,P,I,\theta_0,\Omega,m\rangle
 \longrightarrow_c
 \langle\ell',P',I',\theta_0,\Omega\cup V,m\rangle
}.
\]
Materialization is normalization when instance refinements are already present
in the formula. Assignment evaluates the right side and updates the QCP
storage for the left side. A write through $q$ must update the observation
through $p$ when the state establishes $p=q$. Previous storage contents cannot
remain current facts merely because the pointer identities did not change.

Branches produce separately owned views for $P\land b$ and $P\land\neg b$.
Unknown branches remain; only justified contradictions can prune a branch.
A join either preserves a disjunction or retains common facts. An instance
substitution from one branch cannot be exported to a sibling.

The implementation refreshes at statement and expression boundaries, after
predicate expansion and contract application, and before consuming instance
information. Substitution, merge, and principal write interfaces invalidate
old views. Query-time refresh covers QCP strategies that modify assertions
in place. This is a whole-view invalidation scheme; $\nu$ is not a per-cell
SSA version or a heap address.

\section{Assertion cuts and intermediate instantiation}
For an ordinary assertion $A$, let
\[
 \theta_c=\mathsf{Match}_{QCP}(P,I,A),\qquad
 A^*=\Norm(A[\theta_0][\theta_c]\land\mathsf{Pure}(P)).
\]
The cut emits $P\vdash A^*$ and continues with $(A^*,\Capture(A^*))$. Matching
provides a candidate, not a certificate. A well-scoped substitution for target
existentials satisfies $A^*\vdash A[\theta_0]$ by existential introduction. A pending entailment remains in
$\Omega$ and prevents an unconditional successful result.

The matcher reuses QCP's representation and matching rules for cells, arrays,
fields, spatial predicates, and logical value terms. It specializes each
incoming path separately. A source location reached with a new instance does
not reuse the first visit's candidate substitution.

For example,
\[
 P=x\mapsto3,\quad A=\exists v.\ x\mapsto v\land v\ge0
 \quad\Rightarrow\quad \theta_c=[v\mapsto3],\quad A^*=x\mapsto3.
\]
After an increment, a subsequent cut must carry $4$, not $3$. Pure path facts
survive annotation processing because an annotation does not execute a write.
They refer to logical value identities, with current storage related to those
identities by matching.

If matching cannot recover a relation, the backend retains symbolic terms and
proof obligations. Opaque definitions and abstractions that omit a content
relation can therefore remain precision barriers. A complete per-point view
may contain unknown values; it is not a promise of one closed numeric state.

Which-implies annotations use the same proof-only refinement principle.
After QCP applies the local specification, the current incoming instance
specializes its output. A nontrivial substitution emits an additional
entailment from the original state to that refined output. The original
which-implies obligations remain. This transfer is not applied to mutating
function calls.

In expand mode, \texttt{ignore} leaves proof annotations out of state
abstraction. Under \texttt{check-preserve}, annotations are checked on a copy,
obligations are retained, and the actual execution state continues unchanged.

\section{Function calls}
In no-expand mode, arguments are evaluated in the current product state.
QCP freshens the contract and matches its precondition, obtaining $\theta_g$
and frame $F$. The obligations and successor have the form
\[
 P\vdash R_g[\theta_g]*F,
 \qquad P'=\Norm(E_g[\theta_g]*F),\qquad I'=\Capture(P').
\]
Postcondition existentials are fresh. Old contents in the modified region do
not survive unless the postcondition relates them to the new contents. Frame
facts do survive. A contract with $p\mapsto v$ and $p\mapsto v+1$ transports
input $3$ to output $4$; $\exists w.\ p\mapsto w$ leaves the output symbolic.

Expand mode creates a fresh frame and executes the available body, continuously
updating $(P,I)$. The return transfers the actual heap changes and value, and
removes local storage. A declaration-only call uses its contract and records
that abstraction. A call with neither body nor contract is an error. A
bodyless test target performs precondition checking only; its execution safety
and return-condition phases are not run.

\section{Loops}
Expand mode evaluates the guard and body repeatedly, producing versions
$(P_k,I_k)$ for successive visits. A symbolic guard may fork. Break, continue,
and return transfer their corresponding branch states. Exhausting a bound
reports an incomplete analysis.

No-expand mode has three distinct instance environments:
\begin{align*}
 &P_{\mathrm{entry}}\vdash J[\theta_0],
 &&I_{\mathrm{entry}}=\Capture(P_{\mathrm{entry}}),\\
 &\{J[\theta_0]\land b\}\ \mathit{body}\ \{J[\theta_0]\},
 &&I_{\mathrm{ind}}=\Capture(J[\theta_0]\land b),\\
 &P_{\mathrm{exit}}=\Norm(J[\theta_0]\land\neg b),
 &&I_{\mathrm{exit}}=\Capture(P_{\mathrm{exit}}).
\end{align*}
The induction environment denotes an arbitrary invariant-satisfying iteration.
An entry fact $i=0$ cannot specialize the induction variable to zero. With
$n_0=2$ and invariant $0\le i\le n_0$, the induction range is $0\le i\le2$;
the exit condition $i\ge2$ recovers $i=2$.

Even a fully specified input can yield symbolic intermediate values in this
mode: a loop point summarizes several iterations, and a weak invariant or
contract can admit several outcomes. Preserving instance information means
preserving justified relations, not selecting an arbitrary outcome.

\section{Satisfiability and outcome discipline}
Precondition checking operates on its own normalized copy of $R[\theta_0]$.
It combines accessible predicate definitions, pure conditions, storage ranges,
layout, address bounds, and separation constraints. An UNSAT result for
necessary conditions excludes the original precondition. SAT for necessary
conditions alone does not establish~\eqref{eq:sat}; a complete encoding or a
checked complete witness is required. Residual definitions retain the original
existence obligation for interactive proof.

Initial UNSAT is an invalid instance. Unclosed existence or execution
obligations produce UNKNOWN and residual VCs. Auto/manual classifies whether
the automatic backend closed a VC; replayable certificate status is recorded
separately. A failed entailment after abstraction is not by itself a reachable
execution counterexample.

\section{Relative correctness argument}
\begin{proposition}[Synchronization and execution preservation]
Assume the underlying QCP transitions are sound, normalization preserves
meaning, and substitutions respect types and binder identities. For a valid
synchronized predecessor, a transition satisfying the rules above produces a
successor whose current instance observations are supported by its formula.
When the accumulated obligations hold, concrete executions represented by the
predecessor are covered by the successor representation.
\end{proposition}
\begin{proof}[Argument by transition class]
Ordinary steps inherit QCP's transition property, and capture observes the
result. Normalization preserves denotation. Assertion cuts check the refined
target itself, so the carried relations are justified by the same obligation.
Branches retain separate valuations, and joins use disjunction or common
facts. Calls use the frame rule and fresh post-state variables. Loop induction
uses an independent invariant state; expansion composes individual execution
steps. Budget exhaustion never supplies a success premise.
\end{proof}
This is a relative semantic argument and an implementation contract. The C
modules have executable regression and sanitizer coverage; a machine-checked
proof of their implementation is a separate artifact.

\section{Implementation and reproducibility}
\begin{center}
\begin{tabular}{ll}
\toprule
Component & sac location (under \texttt{SymExec})\\
\midrule
Owned per-assertion view & \texttt{AsrtDef/AssDef.c}\\
Input and current contexts & \texttt{SymExec/InstanceContext.c}\\
Shared normalization & \texttt{Utility/ExpressionNormalize.c}\\
Statement and cut transitions & \texttt{middle\_end/SymExecutor.c}\\
Contract/frame transfer & \texttt{SymExec/FuncCall.c}\\
Optional point observations & \texttt{SymExec/SpecTest.c}\\
Structured output & \texttt{JsonPrinter/JsonFormat.c}\\
\bottomrule
\end{tabular}
\end{center}
Use \texttt{--instance-trace --Json-output} to record statement and condition-expression boundary
formulas, memory observations, equalities, source locations, branches, and
versions. The default artifacts omit the full trace. The new regression suite
is \texttt{SymExec/tests/run-instance-state-tests.py}; it checks positive cases
and stale-value, alias, branch, contract, and induction counterexamples in both
modes. Expected UNKNOWN outcomes establish rejection of an incorrect automatic
success, not a completed proof of the residual goal. Broad example smoke
results are reported separately from semantic expectations.
\end{document}
